\documentclass[11pt,letterpaper]{article}
\usepackage{physical_ai_legislation}
\usepackage[backend=biber,style=numeric,sorting=none]{biblatex}
\addbibresource{fda_national_regulations.bib}

\title{%
  \textbf{FDA Physical AI Oncology Trial Site}\\
  \textbf{National Compliance Guide}\\[6pt]
  \large Federal Regulatory Framework for Physical AI\\
  Oncology Clinical Trial Facilities\\[6pt]
  \normalsize Guidance for integrating Physical AI systems into existing\\
  federal clinical trial, medical device, and drug administration regulations.
}
\author{%
  Kevin Kawchak\\
  CEO, ChemicalQDevice\\
  \texttt{kevink@chemicalqdevice.com}
}
\date{March 24, 2026}

\begin{document}
\maketitle
\thispagestyle{fancy}

\begin{center}
\footnotesize
\textit{This work was adapted from original CFR documents in the public domain.
The original ICH document is copyrighted and may be used, reproduced,
incorporated into other works, adapted, modified, translated or distributed
under a public license. This current work is not endorsed or sponsored by
CFR, ICH, or FDA; and was adapted using Claude Code Opus 4.6.}
\end{center}

\vspace{0.5em}

\begin{center}
\footnotesize
\textit{This draft is independent and is not endorsed, sponsored, or approved
by any trial sponsor, CRO, site, IRB, regulator, or medical society;
and was adapted using Claude Code Opus 4.6.}
\end{center}

\vspace{1em}
\tableofcontents
\newpage

\section{Scope and Purpose}

\subsection{Scope}

This compliance guide identifies the federal regulations applicable to
Physical AI oncology clinical trial sites and specifies where corrections,
additions, and adaptations should be applied to existing regulatory text.
Because a comprehensive amendment of each federal regulation is beyond the
scope of a single document, this guide uses a reference-and-correction approach:
for each applicable regulation, the specific sections requiring updates are
identified, and the nature of the required changes is described.

\subsection{Purpose}

The purpose is to provide Physical AI oncology clinical trial site operators,
sponsors, and regulators with a consolidated roadmap for federal compliance.
This guide is informed by three detailed regulatory adaptations already
completed:

\begin{enumerate}[label=(\arabic*)]
\item Adaptation of ICH E6(R3) for Physical AI oncology clinical trials
  (1,301 lines, 50+ defined terms), covering good clinical practice,
  investigator responsibilities, sponsor responsibilities, and data
  governance~\cite{iche6r3adapt}.
\item Adaptation of 21 CFR Part 50 for Physical AI oncology clinical trials
  (748 lines, 18 Physical AI-specific definitions), covering human subjects
  protection and informed consent~\cite{cfr50adapt}.
\item Adaptation of 21 CFR Part 312 for Physical AI oncology clinical trials
  (1,700 lines, 22 Physical AI definitions), covering IND applications,
  graduated autonomy, safety reporting, and
  decommissioning~\cite{cfr312adapt}.
\end{enumerate}

\subsection{Evidence Base}

The evidence for a Physical AI oncology trial transition is overwhelmingly
based on the most recent simulation~\cite{newtrial2026}, which demonstrated:

\begin{itemize}
\item 168 unique patients across 15 cancer types served in 24 hours by 29
  robot instances of 10 types.
\item Patient throughput 10 times greater than traditional trials (150 to 180
  patients per day versus 5 to 15).
\item Per-patient cost reduction of 75 to 85 percent.
\item Patient retention improvement from 70 to 80 percent (traditional) to
  85 to 95 percent (Physical AI).
\item 99.7 percent facility uptime across 24 continuous hours.
\item Average patient wait time of 8 minutes.
\item 7 adverse events (all Grade 1 or 2, all resolved same hour).
\item Patient satisfaction of 4.7 out of 5.0.
\item Cumulative site PSL of 63.4 to 64.4 (Advanced Site band).
\item Both minute-level clinical detail and 24-hour operational continuity
  handled correctly by AI systems.
\item Staffing reduced from 80 to 120 FTE to 8 to 10 FTE.
\item All clinical procedures, monitoring, documentation, and regulatory
  compliance performed autonomously.
\item Seven surgical procedures with 100 percent R0 resection rate.
\item Peak utilization of 72.4 percent in hour 09 with 28 concurrent patients.
\end{itemize}

\section{21 CFR Part 50 -- Protection of Human Subjects}

\subsection{Sections Requiring Update}

The adapted 21 CFR Part 50 (DOI: 10.5281/zenodo.19040707) provides the
comprehensive text~\cite{cfr50adapt}. The following sections of the original
regulation require corrections:

\begin{enumerate}[label=(\alph*)]
\item \textbf{\S~50.1 (Scope):} Add Physical AI system categories and
  applicability criteria. Specify that the regulation applies to all clinical
  investigations where Physical AI systems perform, assist, or monitor
  procedures. Add minimum USL thresholds by robot type.

\item \textbf{\S~50.3 (Definitions):} Add 18 Physical AI-specific definitions
  including robot agent, USL, digital twin, simulation framework, MCP,
  hash-chained audit trail, HIPAA Safe Harbor de-identification, PCCP,
  federated learning, pre-procedure safety matrix, and conformance levels.

\item \textbf{\S~50.20 (General Requirements for Informed Consent):} Extend
  to require disclosure of robotic system role, autonomy mode, USL score in
  plain language, safety record, emergency stop provisions, digital twin usage,
  and MCP data processing pipeline.

\item \textbf{\S~50.25 (Elements of Informed Consent):} Add Physical AI
  elements as additional required elements, including all 10 robot categories
  and their patient-facing instruction formats~\cite{patientinstructions2026}.

\item \textbf{\S~50.50 through \S~50.56 (Pediatric Protections):} Add
  age-appropriate robotic demonstration requirements, parental consent
  provisions for Physical AI procedures, and advocate appointments.
\end{enumerate}

\subsection{Implementation Notes}

Sites should apply the full adapted text from the Physical AI 21 CFR Part 50
document~\cite{cfr50adapt}. Where the adapted text extends the original
regulation, the extended provisions are additive and do not replace the
original human subjects protection requirements.

\section{21 CFR Part 312 -- Investigational New Drug Application}

\subsection{Sections Requiring Update}

The adapted 21 CFR Part 312 (DOI: 10.5281/zenodo.19057628) provides the
comprehensive text~\cite{cfr312adapt}. Key sections requiring corrections:

\begin{enumerate}[label=(\alph*)]
\item \textbf{\S~312.1-312.2 (Scope):} Extend to cover clinical investigations
  where Physical AI systems perform, assist, or monitor investigational drug
  administration. Define when IND coverage is required based on Physical AI
  risk introduction.

\item \textbf{\S~312.3 (Definitions):} Add 22 Physical AI definitions
  including robot agents, USL scoring (9 platforms evaluated), digital twins,
  simulation frameworks, robot capability profiles, task orders, emergency
  stop, human oversight, MCP, and VLA models.

\item \textbf{\S~312.21 (Phases of Investigation):} Add Phase 0 simulation
  validation requirements. Establish graduated autonomy: Phase 1 teleoperation
  only, Phase 2 graduated autonomy with safety data review, Phase 3 full
  autonomy range with accumulated evidence~\cite{cfr312adapt}.

\item \textbf{\S~312.23 (IND Content):} Add Physical AI System Description
  requirements: system identification, USL score, role in drug administration,
  hardware and software specifications, PCCP, Phase 0 results, digital twin
  specifications, safety architecture, cybersecurity assessment, MCP
  integration, human-robot interaction protocols, and maintenance procedures.

\item \textbf{\S~312.30 (Protocol Amendments):} Define Physical AI-specific
  amendments: new robot type, autonomy level change, significant AI model
  update, simulation framework change, human oversight change, digital twin
  methodology change, and MCP architecture change.

\item \textbf{\S~312.32 (Safety Reporting):} Add seven Physical AI adverse
  event categories: serious Physical AI adverse events, Physical AI-drug
  interactions, cybersecurity incidents (7-day or 15-day reporting), AI model
  performance degradation, sim-to-real divergence (trajectory above 4mm or
  force above 1.0N), digital twin discrepancy, and Physical AI Safety Review
  Committee findings.

\item \textbf{\S~312.33 (Annual Reports):} Add Physical AI elements: USL
  score updates, AI model updates, simulation results, digital twin metrics,
  cybersecurity events, and safety committee findings.

\item \textbf{\S~312.38 (Withdrawal of IND):} Add decommissioning
  requirements: safe deactivation, log preservation, secure patient data
  erasure, 30-day decommissioning timeline, 60-day technical
  analysis~\cite{cfr312adapt}.

\item \textbf{\S~312.40 (Requirements for Use):} Add Physical AI readiness
  requirements: system description submitted, USL threshold met, pre-procedure
  safety matrix verified, training completed, MCP infrastructure configured.

\item \textbf{\S~312.42 (Clinical Holds):} Add eight Physical AI-specific
  grounds: robotic safety failures, AI model degradation, simulation-reality
  divergence (above 2x tolerance), cybersecurity compromise, USL score
  decline, inadequate system description, digital twin failure, and
  inadequate human oversight.
\end{enumerate}

\section{ICH E6(R3) -- Good Clinical Practice}

\subsection{Sections Requiring Update}

The adapted ICH E6(R3) (DOI: 10.5281/zenodo.18973368) provides the
comprehensive text~\cite{iche6r3adapt}. Key sections requiring corrections:

\begin{enumerate}[label=(\alph*)]
\item \textbf{Section 1 (Principles):} Extend foundational principles to
  encompass robotic surgical systems, AI algorithms, and digital twins.
  Add Physical AI system classification (seven robot categories) and the USL
  scoring framework (v1.4.0 through v1.8.0)~\cite{usl2026}.

\item \textbf{Section 2 (Investigator Responsibilities):} Add robotic system
  oversight, AI model monitoring, digital twin validation, and cybersecurity
  incident management to investigator duties.

\item \textbf{Section 3 (Sponsor Responsibilities):} Add quality management
  for Physical AI components: robotic system selection, AI lifecycle
  management, simulation validation, digital twin deployment, and federated
  multi-site coordination.

\item \textbf{Section 4 (Data Governance):} Add full lifecycle governance for
  Physical AI data: robotic telemetry, AI decision logs, digital twin state
  data. Mandate hash-chained audit trails (SHA-256) and HIPAA Safe Harbor
  de-identification. Address five AI categories: generative AI (VLA, diffusion
  policies), agentic AI (LLM-based), reinforcement learning, self-supervised
  learning, and supervised learning.

\item \textbf{Appendix A (Physical AI System Documentation):} Add requirements
  analogous to Investigator's Brochure for Physical AI system descriptions.

\item \textbf{Appendix B (Protocol Templates):} Add Physical AI-specific
  protocol elements including simulation validation, robot deployment plans,
  and graduated autonomy phases.

\item \textbf{Glossary:} Add 50+ terms defined including agentic AI,
  calibration, cobots, digital twins, federated learning, sim-to-real
  transfer, and USL framework.
\end{enumerate}

\section{21 CFR Part 11 -- Electronic Records}

\subsection{Sections Requiring Update}

\begin{enumerate}[label=(\alph*)]
\item \textbf{\S~11.10 (Controls for Closed Systems):} Specify that Physical
  AI system telemetry, AI decision logs, and digital twin records are
  electronic records subject to Part 11 requirements. Hash-chained audit
  trails (SHA-256) satisfy the audit trail requirement when cryptographic
  chain integrity is verified quarterly~\cite{iche6r3adapt}.

\item \textbf{\S~11.30 (Controls for Open Systems):} Apply to federated
  learning data exchanges between Physical AI systems at different sites,
  requiring encryption and digital signature verification.

\item \textbf{\S~11.50 (Signature Manifestations):} Extend to include
  automated system signatures from Physical AI systems that generate,
  modify, or transmit electronic records during clinical procedures.
\end{enumerate}

\section{21 CFR Part 820 -- Quality System Regulation}

\subsection{Sections Requiring Update}

\begin{enumerate}[label=(\alph*)]
\item \textbf{\S~820.30 (Design Controls):} Apply to Physical AI systems
  classified as medical devices, including Phase 0 simulation validation as
  part of design verification and validation.

\item \textbf{\S~820.70 (Production and Process Controls):} Apply to the
  deployment and configuration of Physical AI systems at trial sites,
  including calibration, maintenance, and software update processes.

\item \textbf{\S~820.90 (Nonconforming Product):} Apply to Physical AI
  systems that fall below minimum USL thresholds or exhibit performance
  degradation, with corrective action timelines aligned with the adapted
  21 CFR Part 312 safety reporting requirements~\cite{cfr312adapt}.

\item \textbf{\S~820.198 (Complaint Files):} Extend to include Physical AI
  system malfunction reports, patient complaints involving robotic systems,
  and AI model performance complaints.
\end{enumerate}

\section{Additional Federal Requirements}

\subsection{HIPAA (45 CFR Parts 160 and 164)}

Physical AI oncology clinical trial sites shall comply with HIPAA Privacy,
Security, and Breach Notification Rules with additional attention to the
following Physical AI-specific considerations:

\begin{enumerate}[label=(\alph*)]
\item Robotic telemetry data that can be linked to specific patients
  constitutes protected health information (PHI) under HIPAA.
\item Digital twin data constitutes PHI when it contains or is derived from
  individually identifiable health information.
\item MCP-PAI server architectures shall implement the HIPAA Security Rule
  administrative, physical, and technical safeguards~\cite{iche6r3adapt}.
\item The five MCP-PAI servers (authorization, FHIR, DICOM, ledger,
  provenance) shall each maintain individual HIPAA compliance
  documentation.
\end{enumerate}

\subsection{Nuclear Regulatory Commission (10 CFR Part 35)}

Physical AI oncology clinical trial sites operating radiotherapy systems
shall comply with NRC radiation safety requirements, with the following
Physical AI-specific considerations:

\begin{enumerate}[label=(\alph*)]
\item Radiotherapy positioning robots and motion-tracking robots shall be
  included in the site's radiation safety program as radiation-producing
  equipment accessories.
\item The radiation safety officer shall be trained in Physical AI system
  radiation safety implications.
\item Radiation dose calculations shall account for the precision of Physical
  AI positioning systems, which demonstrated submillimeter accuracy and 94.2
  percent beam gating efficiency in the
  simulation~\cite{newtrial2026}.
\end{enumerate}

\section{Cross-Reference Matrix}

The following table maps Physical AI system categories to the primary federal
regulations that apply to each:

\begin{center}
\footnotesize
\begin{tabular}{lccccc}
\toprule
\textbf{System Type} & \textbf{Part 50} & \textbf{Part 312} &
\textbf{E6(R3)} & \textbf{Part 11} & \textbf{Part 820} \\
\midrule
Surgical robots       & Yes & Yes & Yes & Yes & Yes \\
Cobots                & Yes & Yes & Yes & Yes & Yes \\
RT positioning        & Yes & Yes & Yes & Yes & Yes \\
Needle-placement      & Yes & Yes & Yes & Yes & Yes \\
Social companion      & Yes & --  & Yes & Yes & -- \\
Humanoids             & Yes & --  & Yes & Yes & -- \\
RT motion-tracking    & Yes & Yes & Yes & Yes & Yes \\
Imaging assistants    & Yes & Yes & Yes & Yes & Yes \\
Steerable needles     & Yes & Yes & Yes & Yes & Yes \\
Rehab exoskeletons    & Yes & --  & Yes & Yes & Yes \\
\bottomrule
\end{tabular}
\end{center}

\printbibliography[title={References}]

\end{document}
